WURDE IN 05 ALS SECTION KOPIERT!!!
\section{}{Anforderungsanalyse}
\label{ch:anforderungsanalyse}

Die Anforderungsanalyse bildet die Grundlage für die gesamte Entwicklung
von EffSync. Sie definiert, welche Funktionen das System bereitstellen
muss (funktionale Anforderungen) und welche Qualitätseigenschaften es
erfüllen soll (nicht-funktionale Anforderungen). Die Anforderungen wurden
gemeinsam mit dem Betreuer \textbf{Simon Gunacker} in der Kickoff-Phase
erarbeitet und im Laufe des Projekts iterativ verfeinert.

\subsection{Stakeholder}

Folgende Stakeholder wurden im Rahmen der Anforderungsanalyse identifiziert:

\begin{table}[H]
    \centering
    \caption{Stakeholder-Übersicht}
    \begin{tabularx}{\textwidth}{lXl}
        \toprule
        \textbf{Stakeholder} & \textbf{Interesse / Erwartung}
          & \textbf{Priorität} \\
        \midrule
        Schülerinnen \& Schüler
          & Einfache Aufgabenverwaltung, Kalender, Chat
          & Hoch \\
        Lehrpersonen
          & Überblick über Klassen, Aufgaben und Fortschritt
          & Hoch \\
        Unternehmen (Teams)
          & Skalierbare Projektverwaltung, Sicherheit, DSGVO
          & Hoch \\
        Betreuer (Simon Gunacker)
          & Technische Qualität, Dokumentation, Git-Workflow
          & Mittel \\
        Entwickler (Plavsic, Micail)
          & Wartbarkeit, klare Architektur, Testbarkeit
          & Mittel \\
        \bottomrule
    \end{tabularx}
    \label{tab:stakeholder}
\end{table}

\subsection{Funktionale Anforderungen}
\label{sec:funktionale-anforderungen}

Funktionale Anforderungen beschreiben, \textit{was} das System tun soll.
Sie sind nach Themenbereichen gruppiert und mit einer eindeutigen ID
sowie einer Priorität (Muss / Soll / Kann) versehen.

\subsubsection{}{Benutzerverwaltung \& Authentifizierung}

\begin{table}[H]
    \centering
    \caption{Funktionale Anforderungen -- Benutzerverwaltung}
    \begin{tabularx}{\textwidth}{llXl}
        \toprule
        \textbf{ID} & \textbf{Prio} & \textbf{Anforderung}
          & \textbf{Status} \\
        \midrule
        FA-01 & Muss & Benutzer können sich mit E-Mail und Passwort
                       registrieren und anmelden.
                     & \checkmark \\
        FA-02 & Muss & Passwörter werden gehasht (bcrypt) gespeichert.
                     & \checkmark \\
        FA-03 & Muss & Authentifizierung erfolgt über \ac{JWT}-Token
                       (Access- \& Refresh-Token).
                     & \checkmark \\
        FA-04 & Soll & Anmeldung via Google \ac{OAuth} ist möglich.
                     & \checkmark \\
        FA-05 & Soll & Benutzer können ihr Profil (Name, Avatar,
                       Passwort) bearbeiten.
                     & \checkmark \\
        FA-06 & Kann & Zwei-Faktor-Authentifizierung (\ac{2FA})
                       via \ac{TOTP} (z.\,B.\ Google Authenticator)
                       ist verfügbar.
                     & \checkmark \\
        FA-07 & Muss & Passwort-Reset per E-Mail ist möglich.
                     & \checkmark \\
        \bottomrule
    \end{tabularx}
    \label{tab:fa-auth}
\end{table}

\subsubsection{Rollen \& Berechtigungen}

\begin{table}[H]
    \centering
    \caption{Funktionale Anforderungen -- Rollen \& Berechtigungen}
    \begin{tabularx}{\textwidth}{llXl}
        \toprule
        \textbf{ID} & \textbf{Prio} & \textbf{Anforderung}
          & \textbf{Status} \\
        \midrule
        FA-10 & Muss & Das System verfügt über ein rollenbasiertes
                       Berechtigungssystem (\ac{RBAC}).
                     & \checkmark \\
        FA-11 & Muss & Rollen: Superuser, Admin, Manager, Member,
                       Viewer (Unternehmen) bzw.\ Superuser, Admin,
                       Teacher, Student (Schule).
                     & \checkmark \\
        FA-12 & Muss & Jede Rolle hat klar definierte Berechtigungen
                       für CRUD-Operationen.
                     & \checkmark \\
        FA-13 & Soll & Berechtigungen werden serverseitig geprüft
                       (Guards in NestJS).
                     & \checkmark \\
        FA-14 & Soll & Kritische Operationen werden in einem
                       Audit-Log protokolliert.
                     & \checkmark \\
        \bottomrule
    \end{tabularx}
    \label{tab:fa-rbac}
\end{table}

\subsubsection{Aufgabenverwaltung}

\begin{table}[H]
    \centering
    \caption{Funktionale Anforderungen -- Aufgabenverwaltung}
    \begin{tabularx}{\textwidth}{llXl}
        \toprule
        \textbf{ID} & \textbf{Prio} & \textbf{Anforderung}
          & \textbf{Status} \\
        \midrule
        FA-20 & Muss & Tasks können erstellt, bearbeitet, gelöscht
                       und einem Status zugewiesen werden.
                     & \checkmark \\
        FA-21 & Muss & Tasks haben folgende Felder: Titel,
                       Beschreibung, Status, Priorität,
                       Fälligkeitsdatum, Verantwortliche.
                     & \checkmark \\
        FA-22 & Muss & Tasks können Subtasks enthalten.
                     & \checkmark \\
        FA-23 & Soll & Tasks können Mitgliedern explizit zugewiesen
                       werden (Task-Member-Verwaltung).
                     & \checkmark \\
        FA-24 & Soll & Tasks können nach Status, Priorität und
                       Fälligkeitsdatum gefiltert werden.
                     & \checkmark \\
        FA-25 & Kann & Tasks können mit Tags/Labels versehen werden.
                     & \checkmark \\
        \bottomrule
    \end{tabularx}
    \label{tab:fa-tasks}
\end{table}

\subsubsection{Dashboard \& Metriken}

\begin{table}[H]
    \centering
    \caption{Funktionale Anforderungen -- Dashboard}
    \begin{tabularx}{\textwidth}{llXl}
        \toprule
        \textbf{ID} & \textbf{Prio} & \textbf{An76forderung}
          & \textbf{Status} \\
        \midrule
        FA-30 & Muss & Das Dashboard zeigt eine Übersicht aller
                       aktiven Projekte und Tasks.
                     & \checkmark \\
        FA-31 & Soll & Metriken: Gesamtanzahl Tasks, offene Tasks,
                       abgeschlossene Tasks, überfällige Tasks.
                     & \checkmark \\
        FA-32 & Soll & Statusverteilung wird als Diagramm
                       (Donut-Chart) visualisiert.
                     & \checkmark \\
        FA-33 & Soll & Wöchentliche und monatliche Trends werden
                       als Liniendiagramm dargestellt.
                     & \checkmark \\
        FA-34 & Kann & Fortschrittsbalken für einzelne Projekte
                       werden angezeigt.
                     & \checkmark \\
        FA-35 & Kann & Top-5-Mitglieder nach abgeschlossenen
                       Tasks werden angezeigt.
                     & \checkmark \\
        \bottomrule
    \end{tabularx}
    \label{tab:fa-dashboard}
\end{table}

\subsubsection{Kommunikation \& Benachrichtigungen}

\begin{table}[H]
    \centering
    \caption{Funktionale Anforderungen -- Kommunikation}
    \begin{tabularx}{\textwidth}{llXl}
        \toprule
        \textbf{ID} & \textbf{Prio} & \textbf{Anforderung}
          & \textbf{Status} \\
        \midrule
        FA-40 & Soll & Ein Live-Chat ermöglicht Echtzeit-Kommunikation
                       via WebSockets (Socket.io).
                     & \checkmark \\
        FA-41 & Soll & Chatnachrichten werden persistent gespeichert
                       (DSGVO-konform mit Löschkonzept).
                     & \checkmark \\
        FA-42 & Muss & Benutzer erhalten In-App-Benachrichtigungen
                       bei relevanten Ereignissen (z.\,B.\ neue
                       Task-Zuweisung).
                     & \checkmark \\
        FA-43 & Soll & Benachrichtigungen können als gelesen
                       markiert werden.
                     & \checkmark \\
        FA-44 & Kann & E-Mail-Benachrichtigungen bei kritischen
                       Ereignissen (z.\,B.\ überfällige Tasks).
                     & \checkmark \\
        \bottomrule
    \end{tabularx}
    \label{tab:fa-komm}
\end{table}

\subsubsection{Kalender \& Suche}

\begin{table}[H]
    \centering
    \caption{Funktionale Anforderungen -- Kalender \& Suche}
    \begin{tabularx}{\textwidth}{llXl}
        \toprule
        \textbf{ID} & \textbf{Prio} & \textbf{Anforderung}
          & \textbf{Status} \\
        \midrule
        FA-50 & Soll & Eine Kalenderansicht zeigt Tasks nach
                       Fälligkeitsdatum an.
                     & \checkmark \\
        FA-51 & Soll & Der Kalender unterstützt Monats-, Wochen-
                       und Tagesansicht.
                     & \checkmark \\
        FA-52 & Soll & Eine globale Suche ermöglicht das Auffinden
                       von Tasks, Projekten und Benutzern.
                     & \checkmark \\
        FA-53 & Kann & Suchergebnisse werden in Kategorien
                       gruppiert dargestellt.
                     & \checkmark \\
        \bottomrule
    \end{tabularx}
    \label{tab:fa-kalender}
\end{table}

\subsection{Nicht-funktionale Anforderungen}
\label{sec:nichtfunktionale-anforderungen}

Nicht-funktionale Anforderungen beschreiben \textit{wie} das System
seine Aufgaben erfüllen soll. Sie betreffen Qualitätseigenschaften
wie Performance, Sicherheit und Wartbarkeit.

\begin{table}[H]
    \centering
    \caption{Nicht-funktionale Anforderungen}
    \begin{tabularx}{\textwidth}{llX}
        \toprule
        \textbf{ID} & \textbf{Kategorie} & \textbf{Anforderung} \\
        \midrule
        NFA-01 & Performance
          & API-Antwortzeiten sollen unter 200\,ms liegen
            (bei normaler Last). \\
        NFA-02 & Performance
          & Das Frontend soll innerhalb von 3 Sekunden vollständig
            geladen sein. \\
        NFA-03 & Skalierbarkeit
          & Die Anwendung soll mindestens 100 gleichzeitige
            Benutzer unterstützen. \\
        NFA-04 & Sicherheit
          & Alle Passwörter werden mit bcrypt (Kostenfaktor 12)
            gehasht. \\
        NFA-05 & Sicherheit
          & Die Kommunikation erfolgt ausschließlich über HTTPS. \\
        NFA-06 & Sicherheit
          & \ac{JWT}-Tokens haben eine begrenzte Gültigkeitsdauer
            (Access-Token: 15\,min, Refresh-Token: 7\,Tage). \\
        NFA-07 & Sicherheit
          & Rate Limiting verhindert Brute-Force-Angriffe
            auf Login-Endpunkte. \\
        NFA-08 & Datenschutz
          & Die Anwendung ist \ac{DSGVO}-konform (Löschkonzepte,
            Datensparsamkeit). \\
        NFA-09 & Wartbarkeit
          & Der Code folgt den SOLID-Prinzipien und ist modular
            aufgebaut. \\
        NFA-10 & Wartbarkeit
          & Alle API-Endpunkte sind mit Swagger dokumentiert. \\
        NFA-11 & Usability
          & Die Benutzeroberfläche ist responsiv und auf mobilen
            Geräten nutzbar. \\
        NFA-12 & Usability
          & Die Anwendung unterstützt Dark Mode und Light Mode. \\
        NFA-13 & Verfügbarkeit
          & Das System soll eine Verfügbarkeit von 99\,\% anstreben. \\
        NFA-14 & Testbarkeit
          & Kritische Backend-Module werden durch Unit-Tests
            abgedeckt. \\
        NFA-15 & Kompatibilität
          & Die Anwendung ist mit den aktuellen Versionen von
            Chrome, Firefox und Edge kompatibel. \\
        \bottomrule
    \end{tabularx}
    \label{tab:nfa}
\end{table}

\subsection{Use-Case-Diagramm}
\label{sec:use-cases}

Die folgende Tabelle beschreibt die wichtigsten Use-Cases des Systems
gegliedert nach Akteuren:

\begin{table}[H]
    \centering
    \caption{Use-Case-Übersicht}
    \begin{tabularx}{\textwidth}{llX}
        \toprule
        \textbf{Akteur} & \textbf{UC-ID} & \textbf{Use-Case} \\
        \midrule
        Alle Benutzer
          & UC-01 & Registrieren \& Anmelden \\
          & UC-02 & Profil bearbeiten \\
          & UC-03 & Dashboard einsehen \\
          & UC-04 & Aufgaben anzeigen \& filtern \\
          & UC-05 & Kalender einsehen \\
          & UC-06 & Live-Chat nutzen \\
          & UC-07 & Globale Suche verwenden \\
          & UC-08 & Benachrichtigungen verwalten \\
        \midrule
        Member / Student
          & UC-10 & Aufgaben erstellen \& bearbeiten \\
          & UC-11 & Subtasks verwalten \\
          & UC-12 & Aufgaben als erledigt markieren \\
        \midrule
        Manager / Teacher
          & UC-20 & Projekte \& Klassen verwalten \\
          & UC-21 & Mitglieder zuweisen \\
          & UC-22 & Metriken \& Trends einsehen \\
        \midrule
        Admin
          & UC-30 & Benutzerverwaltung \\
          & UC-31 & Rollen vergeben \\
          & UC-32 & Audit-Log einsehen \\
        \midrule
        Superuser
          & UC-40 & Systemkonfiguration \\
          & UC-41 & Backup \& Restore \\
          & UC-42 & Vollzugriff auf alle Daten \\
        \bottomrule
    \end{tabularx}
    \label{tab:usecases}
\end{table}

\subsection{Detaillierter Use-Case: Aufgabe erstellen}

\begin{table}[H]
    \centering
    \caption{UC-10: Aufgabe erstellen}
    \begin{tabularx}{\textwidth}{lX}
        \toprule
        \textbf{Feld} & \textbf{Beschreibung} \\
        \midrule
        Use-Case-ID   & UC-10 \\
        Name          & Aufgabe erstellen \\
        Akteur        & Member, Manager, Admin, Superuser \\
        Vorbedingung  & Benutzer ist authentifiziert und Mitglied
                        eines Projekts \\
        Nachbedingung & Neue Aufgabe ist in der Datenbank gespeichert
                        und im Frontend sichtbar \\
        Normalablauf  & 1. Benutzer klickt auf \glqq{}Neue Aufgabe\grqq{} \\
                      & 2. Formular mit Titel, Beschreibung, Status,
                           Priorität, Fälligkeitsdatum öffnet sich \\
                      & 3. Benutzer füllt Pflichtfelder aus \\
                      & 4. Benutzer weist optional Mitglieder zu \\
                      & 5. Benutzer speichert die Aufgabe \\
                      & 6. Backend validiert Eingaben und speichert
                           in PostgreSQL \\
                      & 7. Zugewiesene Mitglieder erhalten eine
                           Benachrichtigung \\
        Ausnahmen     & Pflichtfelder fehlen $\rightarrow$ Validierungsfehler \\
                      & Benutzer hat keine Berechtigung $\rightarrow$
                        403 Forbidden \\
        \bottomrule
    \end{tabularx}
    \label{tab:uc10}
\end{table}

\subsection{Detaillierter Use-Case: Anmeldung mit 2FA}

\begin{table}[H]
    \centering
    \caption{UC-01b: Anmeldung mit Zwei-Faktor-Authentifizierung}
    \begin{tabularx}{\textwidth}{lX}
        \toprule
        \textbf{Feld} & \textbf{Beschreibung} \\
        \midrule
        Use-Case-ID   & UC-01b \\
        Name          & Anmeldung mit 2FA \\
        Akteur        & Alle Benutzer (sofern 2FA aktiviert) \\
        Vorbedingung  & Benutzer hat 2FA in den Profileinstellungen
                        aktiviert und eine Authenticator-App
                        eingerichtet \\
        Nachbedingung & Benutzer ist authentifiziert und erhält
                        gültige JWT-Tokens \\
        Normalablauf  & 1. Benutzer gibt E-Mail und Passwort ein \\
                      & 2. Backend prüft Credentials \\
                      & 3. Backend erkennt aktive 2FA und fordert
                           TOTP-Code an \\
                      & 4. Benutzer öffnet Authenticator-App und
                           gibt 6-stelligen Code ein \\
                      & 5. Backend validiert TOTP-Code \\
                      & 6. Access- und Refresh-Token werden ausgestellt \\
        Ausnahmen     & Falscher TOTP-Code $\rightarrow$ 401 Unauthorized \\
                      & Code abgelaufen $\rightarrow$ Neuen Code anfordern \\
        \bottomrule
    \end{tabularx}
    \label{tab:uc01b}
\end{table}